\documentclass[10pt,a4paper]{amsart}
\usepackage[margin=19mm]{geometry}
\usepackage{amsmath,amssymb,mathtools}
\usepackage[scr=rsfs,cal=euler]{mathalfa}
\usepackage{xfrac}
\usepackage[unicode,hidelinks,bookmarks=false]{hyperref}
\newcommand{\ie}{i.e.\ }
\newcommand{\cf}{cf.\ }
\newcommand{\resp}{respectively\ }
\newcommand{\Subsection}{\subsection}
\providecommand{\nobreakdash}{\nobreak\hbox{-}}
\setlength{\emergencystretch}{2em}
\multlinegap0pt
\usepackage{amssymb}
\usepackage{mathtools}
\usepackage[shortlabels]{enumitem}
\usepackage{xcolor}
\newtheorem{proposition}{Proposition}
\title{Port Fillings for Primary Pseudoperfect Numbers}
\author{Han Wang}
\date{Source: arXiv:2605.21518v1 (CC BY 4.0)}
\begin{document}
\maketitle

\noindent\emph{Source and licence.} Port Fillings for Primary Pseudoperfect Numbers. Version arXiv:2605.21518v1 (CC BY 4.0), distributed by arXiv under the Creative Commons Attribution 4.0 International licence, http://creativecommons.org/licenses/by/4.0/, as recorded in the arXiv licence metadata for this exact version and shown on the abstract page https://arxiv.org/abs/2605.21518v1. Full licence notice: \texttt{licenses/arxiv-2605.21518-CC-BY-4.0.txt}.\par\smallskip

\noindent\textit{Editorial scope.} Counted unit: the article's proposition prop:erdos-equivalence (source0.tex lines 138--151). The article's complete original proof of it is delivered with it (source0.tex lines 153--183, one \texttt{proof} environment of 31 lines). No mathematics is written, repaired or re-derived: the statement and the proof are the article's own lines, in the article's own notation, and no retained line is substituted or re-flowed.  No further source-local context is needed: every cross-reference of every retained line resolves inside the two delivered spans.  Nothing was omitted from any retained span, so the package carries no omission record.  No retained line contains a citation, so the wrapper carries no bibliography and no bibliography entry.  No retained line contains a figure environment, an image-inclusion command or drawing code, so no image is delivered and the package declares no asset dependency.  Editorial formatting only, in addition: an AMS \texttt{amsart} preamble that restates the article's own declarations and macros used by the retained lines, namely the environments \texttt{proposition} on separate counters, together with the standard packages those lines need.  The retained environments print in this document as Proposition 1 in order of appearance, whereas the article prints the same items with the numbering of its own full document.  The complete native source of the article as distributed in the arXiv e-print for this exact version is bundled unchanged as \texttt{sources/201067/source0.tex}.
\begin{proposition}\label{prop:erdos-equivalence}
Let $p_1<\cdots<p_k$ be distinct primes and suppose
\[
        \sum_{i=1}^k\frac1{p_i}=1-\frac1m
\]
for some positive integer~$m$.  Then
\[
        m=p_1p_2\cdots p_k.
\]
Consequently the original Erd\H{o}s equation is equivalent to the existence of a primary pseudoperfect number
\[
        N=p_1p_2\cdots p_k.
\]
\end{proposition}

\begin{proof}
Let
\[
        P=p_1p_2\cdots p_k,
        \qquad
        A=\sum_{i=1}^k\frac{P}{p_i}.
\]
Then
\[
        \frac AP=\frac{m-1}{m},
\]
and hence
\begin{equation}\label{eq:mA}
        mA=(m-1)P.
\end{equation}
Fix $i$.  Reducing $A$ modulo~$p_i$, we have
\[
        A\equiv \frac{P}{p_i}\not\equiv0\pmod {p_i},
\]
since all other terms $P/p_j$ with $j\ne i$ are divisible by~$p_i$.  The right-hand side of~\eqref{eq:mA} is divisible by~$p_i$.  Therefore $p_i\mid m$.  This holds for each~$i$, so $P\mid m$.

Write $m=Pt$.  Then
\[
        \frac AP=1-\frac1{Pt},
\]
so
\[
        A=P-\frac1t.
\]
Since $A$ and $P$ are integers, $t=1$.  Thus $m=P$.
\end{proof}

\end{document}
